\documentclass[10pt,a4paper]{amsart}
\usepackage[margin=19mm]{geometry}
\usepackage{amsmath,amssymb,mathtools}
\usepackage[scr=rsfs,cal=euler]{mathalfa}
\usepackage{xfrac}
\usepackage[unicode,hidelinks,bookmarks=false]{hyperref}
\newcommand{\ie}{i.e.\ }
\newcommand{\cf}{cf.\ }
\newcommand{\resp}{respectively\ }
\newcommand{\Subsection}{\subsection}
\providecommand{\nobreakdash}{\nobreak\hbox{-}}
\setlength{\emergencystretch}{2em}
\multlinegap0pt
\usepackage{bm}
\usepackage{bbm}
\usepackage{dsfont}
\usepackage{xspace}
\usepackage{cancel}
\usepackage{mathtools}
\usepackage{amsthm}
\makeatletter
\@ifundefined{theorem}{\newtheorem{theorem}{Theorem}}{}
\makeatother
\newcommand{\Der}{{\operatorname{Der}}}
\newcommand{\Hom}{{\operatorname{Hom}}}
\newcommand{\vphi}{{\varphi}}
\title[Gelfand-Fuks cohomology of vector fields on algebraic variet]{Gelfand-Fuks cohomology of vector fields on algebraic varieties}
\author{Yuly Billig, Kathlyn Dykes}
\date{Source: arXiv:2410.20316v2 (CC BY 4.0)}
\begin{document}
\maketitle

\noindent\emph{Source and licence.} Gelfand-Fuks cohomology of vector fields on algebraic varieties. Version arXiv:2410.20316v2 (CC BY 4.0), distributed under the Creative Commons Attribution 4.0 International licence, as recorded in the arXiv licence metadata for this exact version and shown on the abstract page https://arxiv.org/abs/2410.20316v2. Full rights notice: licenses/arxiv-2410.20316-CC-BY-4.0.txt.\par\smallskip

\noindent\textit{Editorial scope.} Counted unit: the Theorem whose source label is \texttt{theorem:Kunnethsemidirect} in the article (source0.tex lines 1257--1271, source label \texttt{theorem:Kunnethsemidirect}), delivered together with the article's own complete original proof of it (source0.tex lines 1352--1403, one \texttt{proof} environment of 52 lines). No mathematics is written, repaired or re-derived: statement and proof are the article's own text line for line. The proof is detached from the statement in the source: the article states the result at lines 1257--1271 and prints its proof at lines 1352--1403, and the proof environment's own header names the statement, so the association is the article's own. Required context retained uncounted: none. Every definition, display and label of those lines is retained unchanged. Omitted from this package and not redistributed: 1--1256, 1272--1351, 1404--1558. Nothing inside a retained excerpt is omitted or altered: every retained body is delivered exactly as the article's lines are written, so no omission receipt of any kind accompanies this package. Citations: the retained excerpts carry no citation command at all, so no bibliography entry of the article is transcribed and no reference is redistributed. Reference adaptation: none was needed and none was made; every reference inside the delivered unit resolves inside it, so no unresolved reference or citation is produced. Printed slips kept literal: no printed word, symbol, hypothesis or number of the counted statement or of its proof was altered. Layout and class: the article's own class is \texttt{amsart} with packages including geometry, ifdraft, amsmath, amsthm, amssymb, amsfonts, enumitem, hyperref, graphicx, mathtools, mathrsfs, bbm, stmaryrd, tikz-cd; the wrapper is the lane's pinned \texttt{amsart} editorial wrapper at 10pt on A4, which restates the theorem-like environments the retained text needs and adds the article's own macro definitions that the retained text needs (\texttt{\textbackslash Der}, \texttt{\textbackslash Hom}, \texttt{\textbackslash vphi}), with extra wrapper packages (bm, bbm, dsfont, xspace, cancel, mathtools). The delivered numbering is the unit's own. Assets: no retained span contains a figure environment, a graphics-inclusion command, a PDF-page inclusion, a table or a TikZ picture; no image file is copied into this package, the package declares no asset dependency and no image file is redistributed with it. Licence: version arXiv:2410.20316v2 of this article is distributed under the Creative Commons Attribution 4.0 International licence, as recorded in the arXiv licence metadata for this exact version and shown on the abstract page https://arxiv.org/abs/2410.20316v2. A full rights notice is delivered at licenses/arxiv-2410.20316-CC-BY-4.0.txt. Characters: every retained excerpt is pure ASCII, so the delivered engine, XeLaTeX with its default Latin Modern text font, reports no missing character.
\begin{theorem}\label{theorem:Kunnethsemidirect}
Let $(A, V)$ be a Lie Rinehart pair, where the action of $V$ acting on $A$ is given by $\rho: V \rightarrow \Der(A)$. Let $L$ be a Lie algebra and $W$ an $L$-module. Consider the semi-direct product $V \ltimes (A \otimes L)$ of Lie algebras %with $\rho: V \rightarrow \Der(A)$ 
such that
\[
[v_1 + f_1 \otimes x_1, v_2 + f_2 \otimes x_2] = [v_1, v_2] + f_1f_2 \otimes[x_1, x_2] + \rho(v_1)f_1 \otimes x_1 - \rho(v_2)f_2 \otimes x_2
\]
for $v_1, v_2 \in V$, $f_1, f_2 \in A$ and $x_1, x_2 \in L$. For the $V \ltimes (A \otimes L)$-module $A \otimes W$  with action 
\[
(v + f \otimes x) (g \otimes w) = \rho(v)(g) \otimes w + fg \otimes x(w)
\]
for $f,g \in A, v \in V, w \in W$, then the following K\"unneth formula holds:
\[
H_A^{n}(V\ltimes (A \otimes L), A \otimes W) \cong \sum_{i+j =n} H_A^i(V, A) \otimes H^j(L, W) .
\]
\end{theorem}

\begin{proof}[Proof of Theorem \ref{theorem:Kunnethsemidirect}]
The restriction of the map $\ast$ to cocycles will have codomain $Z_A(V\ltimes (A \otimes L), A \otimes W)$ as $d(\alpha \ast \beta) =0$ if both $d\alpha =0$ and $d\beta =0$. To induce a map on cohomologies, we need the images of $B_A^i(V, A) \otimes Z^{n-i}(L,W)$ and $Z_A^i(V,A) \otimes B^{n-i}(L,W)$ to be contained in $B_A^n(V \ltimes (A \otimes L), A \otimes W)$. Suppose $\alpha \in B_A^i(V,A)$ then $\alpha = d\delta$ for some $\delta \in C_A^{n-i}(V, A)$. Then $d(\delta \ast \beta) = d\delta \ast \beta= \alpha \ast \beta$ for every $\beta \in Z^{n-i}(L, W)$. Similarly, $Z_A^i(V, A) \ast B^{n-i}(L,W) \in B^n_A(V \ltimes(A \otimes L), A \otimes W)$.  This induces a map on homologies:
\[
\ast: \sum_{i=0}^n H_A^i(V, A) \otimes H^{n-i}(L, W) \rightarrow H_A^{n}(V \ltimes (A \otimes L), A \otimes W).
\]
We will prove this map is bijective. 

First remark that as vector spaces, $V \ltimes (A \otimes L) \cong V \oplus (A \otimes L)$ so that 
\begin{equation}\label{equation:vectorspdecomp}
\bigwedge_A^n(V \ltimes (A \otimes L)) \cong \bigoplus_{i=0}^n \bigwedge_A^i V \wedge \bigwedge_A^{n-i} (A \otimes L)
\end{equation}
as this is a product over $A$. Let $\vphi \in C^n_A(V \ltimes (A \otimes L), A \otimes W)$ be arbitrary. Then $\vphi$ decomposes into $\vphi = \sum_{i=0}^n \delta_i \wedge \gamma_{n-i}$ for  $\delta_i \in C_A^i(V, A)$ and $\gamma_{n-i} \in C_A^{n-i}(A \otimes L, W)$.

For injectivity, suppose that $\sum_{i=0}^n \alpha_i \ast \beta_{n-i} \in B(V \ltimes (A \otimes L), A \otimes W)$ for some $\alpha_i \in C^i_A(V, A)$ and $\beta_j \in C^j(L, W)$. Then there exists $\omega \in C^{n-1}_A(V \ltimes (A \otimes L), A \otimes W)$ such that $\sum_{i=0}^n \alpha_i \ast \beta_{n-i} = d \omega$. By the above remark, $\omega$ decomposes as  $\sum_{i=0}^{n-1} \delta_i \wedge \gamma_{n-1-i}$ for $\delta_i \in C_A^i(V, A)$ and $\gamma_{j}\in C_A^j(A\otimes L, W)$. Then
\begin{align*}
\sum_{i=0}^n \alpha_i \ast \beta_{n-i} & = \delta_0 \wedge d\gamma_{n-1} + \sum_{i=1}^{n-1} \left( \delta_{i} \wedge d\gamma_{n-i-1} +  d\delta_{i-1} \wedge \gamma_{n-i} \right) + d\delta_{n-1} \wedge \gamma_0 
\end{align*}
By restricting to each term in the decomposition (\ref{equation:vectorspdecomp}),  $\alpha_0 \ast \beta_{n} = \delta_0 \wedge d\gamma_{n-1}$,  $\alpha_n \ast \beta_0 = d\delta_{n-1} \wedge \gamma_0$ and  $\alpha_i \ast \beta_{n-i} = \delta_{i} \wedge d\gamma_{n-i-1} +  d\delta_{i-1} \wedge \gamma_{n-i}$ for $1 \leq i \leq n-1$. 

After restricting to $\bigwedge^{n} 1 \otimes L$, the first equality gives $\alpha_0 = \delta_0$ and $\beta_n = d\gamma_{n-1}$. We now proceed by induction on $i$. Suppose that $\beta_{n-(i-1)} \in B^{n-(i-1)}(L,W)$ and that $\alpha_{i-1} = \delta_{i-1}$. Then $d \alpha_{i-1} =0 \implies d\delta_{i-1}=0$ so that the relation reduces to 
\[
\alpha_i \ast \beta_{n-i} = \delta_{i} \wedge d\gamma_{n-i-1} +  d\delta_{i-1} \wedge \gamma_{n-i} = \delta_i \wedge d \gamma_{n-i-1}
\]
so that $\alpha_i = \delta_i$ and $\beta_{n-i} \in B^{n-i}(L,W)$ for all $0 \leq i \leq n-1$. Finally, when $i =n$, we have $\alpha_n \ast \beta_0 = d \delta_{n-1} \wedge \gamma_0$, but since $\delta_{n-1} = \alpha_{n-1} \in Z_A(V, A)$ so that $\alpha_n \ast \beta_0 =0$, thus either $\alpha_n \in B^n_A(V, A)$ or $\beta_0 \in B^0(L, W)$. Thus $\alpha_i \otimes \beta_{n-i} \in Z_A^i(V,A) \otimes B^{n-i}(L,W)$ and $\ast$ is injective. 

To show surjectivity, we will show for every $\vphi \in Z_A^{n}(V \ltimes (A \otimes L), A \otimes W)$, there exists some $\alpha_i \in Z_A^i(V, A)$ and $ \beta_{n-i} \in Z^{n-i}(L, W)$ such that $\vphi = \sum_{i=0}^n \alpha_i \ast \beta_{n-i}$. 

By our remark, $\vphi = \sum_{i=0}^n \delta_i \wedge \gamma_{n-i}$ for some appropriate $\delta$ and $\gamma$. Set $\alpha_i(v_1, \dots, v_i) = \delta_i(v_1, \dots, v_i)\in C^i_A(V,A)$ and $\beta_j (x_1, \dots, x_j) = \gamma_j(1\otimes x_1, \dots, 1\otimes x_j)\in C^j(L, W)$. As the action of $A$ commutes with $L$, we see that 
\begin{align*}
\alpha_i \ast \beta_{n-i} (v_1, \dots, v_i, f_1 \otimes x_1, \dots, f_{n-i} \otimes x_{n-i}) & = f_1 \cdots f_{n-i}\alpha(v_1, \dots, v_i) \otimes \beta(x_1, \dots, x_{n-1})\\
& = f_1 \cdots f_{n-i} \delta_i(v_1, \dots, v_i) \otimes \gamma_{n-i}(1 \otimes x_1, \dots, 1 \otimes x_{n-1})\\
& = \delta_i(v_1, \dots, v_i) \otimes \gamma_{n-i}(f_1 \otimes x_1, \dots, f_{n-i} \otimes x_{n-i}).
\end{align*}
Thus $\vphi = \sum_{i=0}^n \alpha_i \ast \beta_{n-i}$. 


Now, consider $1 \leq i \leq n$. We have that $d_A\alpha_{i-1} \ast \beta_{n+1-i} + (-1)^i\alpha_{i} \ast d\beta_{n-i}=0$. We will show this can only hold if both terms are themselves zero. 

If either $\beta_{n+1-i}$ or $\alpha_i$ are zero, then we can assume that $\alpha_{i-1}$ or $\beta_{n-i}$ are zero respectively. We would conclude that $\beta_{n-i}$ or $\alpha_{i-1}$ is a cocycle.  Suppose that neither derivations are zero. Then $\alpha_i = d_A\alpha_i$ and $\beta_{n+1-i} = d \beta_{n-i}$. 
\[
d(\alpha_{i-1} \ast \beta_{n-i}) = \alpha_{i-1} \ast d \beta_{n-i} + d\alpha_{i-1} \ast \beta_{n-i} = \alpha_{i-1} \ast \beta_{n+1-i} + \alpha_i \ast \beta_{n-i}
\]
so that this part of $\vphi$ is a coboundary and we can remove it. After removing all such coboundaries, we would have that every $\alpha_i$, $\beta_i$ is a cocycle. 

\iffalse
Consider $i=1$. Then $d_A \alpha_0 \ast \beta_n = \alpha_1 \ast d \beta_{n-1}$. We know  $d_A \alpha_0 (v) = v(a)$ for any $v \in V$. Then $d_A\alpha_0 \ast \beta_n (v, f_1 \otimes x_1, \dots, f_{n-1} \otimes x_{n-1}) = f_1 \cdots f_{n-1} v(a) \otimes \beta_n(x_1, \dots, x_{n-1})$. 

For $1 \leq i \leq n$,  By evaluating at $(v_1, \dots, v_i, 1, \dots, 1)$ we see that $d_A\alpha_{i-1} = (-1)^i \alpha_i$ and by evaluating at $(1, \dots, 1, 1 \otimes x_1, \dots, 1 \otimes x_{n-i})$ we see that $\beta_{n+1-i} = d\beta_{n-i}$. Thus for $1 \leq i \leq n$, $\alpha_i \in B^i_A(V,A)$ and for $1 \leq i \leq n$, we have that $\beta_i \in B^i (L,W)$. We need to show that $\alpha_0$ and $\beta_0$ are also cocycles. 

First, note that $C^0_A(V,A) = \Hom_A (\bigwedge^0 V, A) = \Hom_A(A,A) \cong A$ so that $\alpha_0 = a_0$ for some $a_0 \in A$. Then $a_0 \ast d \beta_n(f_1 \otimes x_1, \dots, f_n \otimes x_n) = f_1\cdots f_n a_0 \otimes d\beta_n(x_1, \dots, x_n) = 1 \otimes f_1 \cdots f_n a_0 d\beta_n(x_1, \dots, x_n)$ so that either $a_0=0$ or $d\beta_n =0$. 
\fi
\end{proof}

\end{document}
